\documentclass[12pt]{article}
\usepackage[utf8]{inputenc}
\usepackage[T1]{fontenc}
\usepackage[margin=1in]{geometry}
\usepackage{setspace}
\usepackage{xcolor}
\usepackage[colorlinks=true, urlcolor=blue!60!black, linkcolor=blue!60!black]{hyperref}
\setlength{\parindent}{2em}
\setlength{\parskip}{1.25em}
\renewcommand{\baselinestretch}{1.0}
\begin{document}
\begin{center}
{\LARGE \textit{NFIB v. Sebelius}, 567 U.S. 519 (2012)\par}
\vspace{0.8em}
{\large Case Brief --- Four-Page Edition\par}
\end{center}
\vspace{1em}

\section{Memory Jogger}

Chief Justice Roberts joined the four liberal Justices to uphold the Affordable Care Act's individual mandate as a tax---even though five Justices agreed it exceeded the Commerce Clause---and joined six other Justices to hold that threatening States with the loss of all existing Medicaid funding to force the Medicaid expansion was unconstitutionally coercive (\href{https://supreme.justia.com/cases/federal/us/567/519/}{Justia}).

\section{Detailed Case Facts}

The 2010 Affordable Care Act (ACA) sought to expand coverage in a \$2.5 trillion health market in which roughly 50 million residents were uninsured and hospitals had to treat the uninsured, shifting costs to insured families. Its ``guaranteed-issue'' and ``community-rating'' reforms barred insurers from denying coverage or raising premiums for health status, which invites healthy people to wait until they are sick to buy insurance. Congress paired them with two challenged provisions. \textbf{The individual mandate} required most Americans to maintain ``minimum essential'' coverage or pay a ``shared responsibility payment,'' labeled a ``penalty,'' to the IRS with their tax returns; it was capped at the price of a bronze plan, exempted those below the filing threshold, and could not be enforced by prosecution, liens, or levies. \textbf{The Medicaid expansion} required States to cover nearly all adults up to 133\% of the poverty line, with the federal government paying 100\% of new costs at first and 90\% later; a State that refused could lose \textit{all} its federal Medicaid funds, which exceeded 10\% of most state budgets.

\section{Procedural History}

Twenty-six States, individuals, and NFIB sued. The district court invalidated the whole ACA. The Eleventh Circuit held the mandate beyond the Commerce Clause but severable, and upheld the Medicaid expansion. Amid a circuit split, the Court granted certiorari, heard three days of argument in March 2012, and decided the case on June 28, 2012.

\section{Judicial Votes}

Affirmed in part and reversed in part, by shifting coalitions.
\begin{itemize}
\item \textbf{Anti-Injunction Act (9--0):} no bar to pre-enforcement review.
\item \textbf{Commerce and Necessary and Proper Clauses (5--4 against the mandate):} Roberts, writing for himself, and the joint dissenters (Scalia, Kennedy, Thomas, Alito); Ginsburg, Breyer, Sotomayor, and Kagan would have upheld it.
\item \textbf{Taxing Power (5--4 upholding the mandate):} Roberts for the Court, joined by Ginsburg, Breyer, Sotomayor, and Kagan.
\item \textbf{Medicaid coercion (7--2):} Roberts, Breyer, and Kagan (plurality) plus the joint dissenters; Ginsburg and Sotomayor dissented.
\item \textbf{Remedy (5--4):} Roberts, Breyer, and Kagan, with Ginsburg and Sotomayor, severed only the threat to existing funds; the joint dissenters would have voided the entire Act. Thomas also dissented separately.
\end{itemize}

\section{Holding}

Because Congress called the payment a ``penalty,'' the Anti-Injunction Act does not bar suit. Looking past labels, an exaction that is not prohibitory, lacks scienter, is collected by the IRS without criminal sanctions, raises revenue, and leaves a lawful choice to pay is a valid tax, not a direct tax requiring apportionment. Congress coerces States unconstitutionally when it threatens an independent, existing grant amounting to a large share of state budgets to compel acceptance of a new program effecting a ``shift in kind, not merely degree.'' Whether Roberts's Commerce Clause analysis is holding or dicta remains debated.

\section{Analysis of Opinions}

\textbf{Roberts.} The federal government has only enumerated powers. The power to ``regulate Commerce'' presupposes existing activity; compelling people to buy an unwanted product---like ordering them to buy broccoli---would let Congress regulate citizens ``from cradle to grave.'' Even \textit{Wickard}'s farmer was growing wheat. Nor was the mandate ``proper,'' since it created the very commerce it regulated. But because it was ``fairly possible'' to read \S~5000A as a tax, and it met \textit{Bailey v. Drexel Furniture}'s functional test, it was valid; the taxing power cannot criminalize inactivity, only require payment. On Medicaid, whereas \textit{South Dakota v. Dole} threatened under 0.5\% of a State's budget, the ACA threatened over 10\%, ``economic dragooning'' toward a transformed program of which States had no fair notice. The remedy was to bar withdrawal of existing funds.

\textbf{Ginsburg (concurring in part and dissenting in part).} The activity/inactivity line is a ``crabbed reading'' recalling pre-1937 formalism; nearly everyone enters the health-care market, and self-insuring is an economic decision. Under \textit{Gonzales v. Raich} the mandate was necessary and proper to prevent an adverse-selection ``death spiral.'' The Medicaid expansion amended an existing program, and the coercion test lacks manageable standards.

\textbf{Joint dissent (Scalia, Kennedy, Thomas, Alito).} Treating a ``penalty'' as a tax was ``verbal wizardry''; the mandate and Medicaid expansion were the Act's pillars, so the whole Act must fall. \textbf{Thomas} added that the ``substantial effects'' test is inconsistent with the original understanding.

\section{Examples: Future Applications}

\begin{enumerate}
\item \textbf{Same-Side (Valid Tax): Flood-Insurance Assessment.} A modest IRS-collected \$400 assessment on flood-zone homeowners who lack flood insurance, with no criminal sanction or scienter, satisfies Part III--C.
\item \textbf{Same-Side (Coercive): Pre-K Ultimatum.} Withholding all existing Title I funds unless a State adopts a new universal preschool program couples a ``shift in kind'' with a threat to entrenched funds.
\item \textbf{Opposite-Side (Invalid): Electric-Vehicle Mandate.} Ordering every driver to buy an electric vehicle on pain of a \$15,000 fine and jail exceeds the commerce power and cannot be saved as a tax.
\item \textbf{Opposite-Side (Valid Condition): 4\% of Highway Funds.} Conditioning a small share of highway grants on work-zone speed cameras mirrors \textit{Dole}.
\item \textbf{Fence-Sitter: Steep or Zero Exactions.} A \$2,500 IRS assessment for skipping a \$50 vaccine may be so prohibitory that it becomes a penalty; a \$0 exaction, as in \textit{California v. Texas}, raises no revenue at all.
\end{enumerate}

\section{Critique}

\textbf{Originalist and textualist views.} Randy Barnett hailed the Commerce Clause ruling as shifting the ``constitutional gestalt'' toward enforceable limits, while then-Professor Amy Coney Barrett argued the tax construction pushed the text ``beyond its plausible meaning.'' Michael Greve welcomed the coercion holding but found its budget-percentage formula unresolved.

\textbf{Progressive views.} Critics call the line between activity and inactivity arbitrary and reminiscent of \textit{Lochner}, question the originalist pedigree of strict enumerationism, and argue that ``active avoidance'' let the Court announce new limits without paying the cost of striking the law, while the Medicaid holding created a coverage gap for millions.

\section{Key Quotations}

\begin{quote}
``The Framers gave Congress the power to regulate commerce, not to compel it, and for over 200 years both our decisions and Congress's actions have reflected this understanding. There is no reason to depart from that understanding now.''
\end{quote}
\noindent --- Chief Justice Roberts, 567 U.S. at 555.

\begin{quote}
``In this case, the financial `inducement' Congress has chosen is much more than `relatively mild encouragement'---it is a gun to the head. \ldots\ The threatened loss of over 10 percent of a State's overall budget, in contrast, is economic dragooning that leaves the States with no real option but to acquiesce in the Medicaid expansion.''
\end{quote}
\noindent --- Chief Justice Roberts, plurality opinion, 567 U.S. at 581--82.

\begin{quote}
``The Affordable Care Act's requirement that certain individuals pay a financial penalty for not obtaining health insurance may reasonably be characterized as a tax. Because the Constitution permits such a tax, it is not our role to forbid it, or to pass upon its wisdom or fairness.''
\end{quote}
\noindent --- Chief Justice Roberts, for the Court, 567 U.S. at 574.

\end{document}
